\section{Methods}
\label{sec:methods}     % 内部引用符，文章中其他部分可以引用这里

% 本文从监督目标、推理接口和 rationale 来源三个方面考察模型性能。
% 仅监督最终分数的训练方式记为 score-only supervision（SO），
% 同时监督 rationale 和最终分数的训练方式记为
% rationale-and-score supervision（RS）。直接输出分数的推理接口
% 记为 direct-score inference（DS），先生成 rationale 再输出分数
% 的推理接口记为 rationale-first inference（RF）。以标准 RS 为基线，
% 我们进一步比较 ReaSon、PaIR 和 RePaIR 三种方法。

We study model performance along three axes: supervision target,
inference protocol, and rationale source. Training that supervises
only the final score is referred to as score-only supervision (SO),
whereas training that supervises both the rationale and final score
is referred to as rationale-and-score supervision (RS). Direct-score
inference (DS) outputs the score directly, while rationale-first
inference (RF) generates a rationale before the score. Using standard
RS as the baseline, we further compare three methods: ReaSon, PaIR,
and RePaIR.

% 下面首先形式化标准 RS。设 \(x_i\) 为第 \(i\) 个样本的输入，
% \(s_i\) 为真实分数，\(r_i^T\) 为教师模型生成的 rationale。
% 其监督目标为
% \[
% \begin{aligned}
% y_i^R ={}&
% \texttt{\textless reasoning\textgreater}\,r_i^T\,
% \texttt{\textless/reasoning\textgreater} \\
% &\texttt{\textless score\textgreater}\,s_i\,
% \texttt{\textless/score\textgreater}.
% \end{aligned}
% \]
% 标准 RS 对整个 completion 中的 token loss 统一取平均。
% 因此，其训练目标可以分解为
% \[
% \begin{aligned}
% \mathcal{L}_{\mathrm{RS}}
% ={}&
% \frac{N_r}{N_r+N_s}\mathcal{L}_r \\
% &+
% \frac{N_s}{N_r+N_s}\mathcal{L}_s ,
% \end{aligned}
% \]
% 其中，\(N_r\) 和 \(N_s\) 分别表示 rationale 和 score 部分
% 的 token 数量，\(\mathcal{L}_r\) 和 \(\mathcal{L}_s\)
% 表示对应的平均 token-level loss。当 \(N_r \gg N_s\) 时，
% rationale 部分在总体标量损失中具有更大的权重。

We first formalize standard RS. Let \(x_i\) denote the input of the
\(i\)-th example, \(s_i\) its gold score, and \(r_i^T\) the rationale
generated by the teacher model. The target completion is
\[
\begin{aligned}
y_i^R ={}&
\texttt{\textless reasoning\textgreater}\,r_i^T\,
\texttt{\textless/reasoning\textgreater} \\
&\texttt{\textless score\textgreater}\,s_i\,
\texttt{\textless/score\textgreater}.
\end{aligned}
\]
Standard RS averages token losses over the entire completion.
Its objective can therefore be decomposed as
\[
\begin{aligned}
\mathcal{L}_{\mathrm{RS}}
={}&
\frac{N_r}{N_r+N_s}\mathcal{L}_r
&+
\frac{N_s}{N_r+N_s}\mathcal{L}_s ,
\end{aligned}
\]
where \(N_r\) and \(N_s\) are the numbers of rationale and score
tokens, and \(\mathcal{L}_r\) and \(\mathcal{L}_s\) are their
respective mean token-level losses. When \(N_r \gg N_s\), the
rationale span receives a larger weight in the overall scalar loss.

% \subsection{ReaSon: Regenerated Rationales with Gold-Score Restoration}

% ReaSon 在不改变标准 RS 训练目标的情况下，将教师 rationale
% 替换为学生模型自生成的 rationale。这里的 regeneration
% 仅指重新生成 rationale 并恢复 gold score，不包含显式的
% critique--revision loop。
% 教师 rationale \(r_i^T\) 由 \texttt{deepseek-v4-pro} 蒸馏，
% 经过格式和 gold-score 一致性筛选后构成
% \(\mathcal{D}_T=\{(x_i,r_i^T,s_i)\}\)。第一阶段在
% \(\mathcal{D}_T\) 上训练模型 \(\theta_A\)，再由该模型通过
% RF 生成新的 rationale \(r_i^A\)。生成分数仅用于记录，
% 训练时仍使用 gold score，由此得到
% \(\mathcal{D}_A=\{(x_i,r_i^A,s_i)\}\)。
% 第二阶段从相同的基础模型重新训练，而不是继续更新
% \(\theta_A\)。因此，ReaSon 的主要变化是将训练 rationale
% 从 \(r_i^T\) 替换为 \(r_i^A\)。最终模型在 DS 和 RF
% 接口下分别评估。数据构造见附录 A，prompt 和生成配置见附录 B。

\subsection{ReaSon: Regenerated Rationales with Gold-Score Restoration}

ReaSon replaces teacher rationales with student-generated rationales
while retaining the standard RS objective. Here, regeneration refers
only to rationale regeneration followed by gold-score restoration.
It does not involve an explicit critique--revision loop.

Teacher rationales \(r_i^T\) are distilled from
\texttt{deepseek-v4-pro} and filtered for format validity and
gold-score consistency, forming
\(\mathcal{D}_T=\{(x_i,r_i^T,s_i)\}\). A first-stage model
\(\theta_A\) is trained on \(\mathcal{D}_T\) and then generates new
rationales \(r_i^A\) under RF. Its predicted scores are retained only
for recording, while the gold scores remain the training labels.
This produces \(\mathcal{D}_A=\{(x_i,r_i^A,s_i)\}\).

The second-stage model is reinitialized from the same base model
rather than obtained by further updating \(\theta_A\). ReaSon
therefore primarily replaces \(r_i^T\) with \(r_i^A\) as the source
of rationale supervision. The resulting model is evaluated under
both DS and RF. Data construction is described in Appendix A, while
prompt and generation details are provided in Appendix B.

% \subsection{Paired-Interface Reweighting (PaIR)}

% 标准 RS 存在两个局限：较长的 rationale 在 token 平均损失中
% 具有更大权重，而且训练过程只覆盖 RF 接口。为缓解前者，
% 我们在单条 RF 序列内分别加权 rationale 和 score 区域，
% 并阻断 score 对 rationale 的注意力。然而，该设计仍然只使用
% RF prompt，因而无法解决 DS 接口失配。

% PaIR 为每个样本构造配对的 DS 和 RF 训练视图。两个视图共享
% 相同的 user prompt 和 gold score，但采用不同的 system prompt
% 和输出结构。DS 目标 \(y_i^D\) 仅包含 score block
% \(\texttt{\textless score\textgreater}\,s_i\,
% \texttt{\textless/score\textgreater}\)。RF 目标 \(y_i^R\)
% 则在相同的 score block 之前增加 rationale block
% \(\texttt{\textless reasoning\textgreater}\,r_i\,
% \texttt{\textless/reasoning\textgreater}\)。

% 两个视图的平均 token-level loss 分别记为
% \(\mathcal{L}_D\) 和 \(\mathcal{L}_R\)，并进行等权组合：
% \[
% \mathcal{L}_{\mathrm{PaIR}}
% =0.5\mathcal{L}_D+0.5\mathcal{L}_R.
% \]
% 组合后的标量损失仅执行一次反向传播和参数更新。
% 该目标平衡的是 DS 和 RF 两个视图，而不是 RF 视图内部的
% rationale 和 score token。

% 我们以 Mix 作为消融对照。Mix 与 PaIR 使用相同的数据和训练配置，
% 仅取消分支级归一化。PaIR 在 DS 和 RF 上均优于 Mix，
% 表明分支级归一化能够带来额外收益。

\subsection{Paired-Interface Reweighting (PaIR)}

Standard RS has two limitations: long rationales receive greater
weight under token-averaged training, and only the RF interface is
observed during training. To address the former, we separately
weighted the rationale and score regions within a single RF sequence
and prevented the score region from attending to the rationale.
However, this design still uses only RF prompts and therefore cannot
resolve the DS interface mismatch.

PaIR constructs paired DS and RF training views for each example.
The two views share the same user prompt and gold score but differ in
their system prompts and output structures. The DS target \(y_i^D\)
contains only the score block
\(\texttt{\textless score\textgreater}\,s_i\,
\texttt{\textless/score\textgreater}\). The RF target \(y_i^R\)
adds a rationale block
\(\texttt{\textless reasoning\textgreater}\,r_i\,
\texttt{\textless/reasoning\textgreater}\) before the same score
block.

The mean token-level losses of the two views are denoted by
\(\mathcal{L}_D\) and \(\mathcal{L}_R\). They are computed separately
and combined with equal weight:
\[
\mathcal{L}_{\mathrm{PaIR}}
=0.5\mathcal{L}_D+0.5\mathcal{L}_R.
\]
The resulting scalar loss is used for a single backward pass and
parameter update. This objective balances the DS and RF views rather
than the rationale and score tokens within the RF view.

We use Mix as a matched ablation. It shares the same data and training
configuration as PaIR but removes branch-wise normalization. PaIR
outperforms Mix under both DS and RF, indicating that branch-wise
normalization provides additional gains.

% \subsubsection{Regenerated-Rationale Paired-Interface Reweighting (RePaIR)}

% RePaIR 将 ReaSon 的 rationale 替换策略与 PaIR 的配对监督
% 方式相结合。第一阶段从基础模型 \(\theta_0\) 出发，在教师
% rationale 数据集 \(\mathcal{D}_T\) 上训练 PaIR 模型
% \(\theta_A\)。随后，\(\theta_A\) 在 RF 接口下重新生成
% rationale \(r_i^A\)，并与 gold score \(s_i\) 组合为
% \(\mathcal{D}_A=\{(x_i,r_i^A,s_i)\}\)。

% 第二阶段重新从 \(\theta_0\) 初始化模型，并在
% \(\mathcal{D}_A\) 上进行 PaIR 训练，而不是继续更新
% \(\theta_A\)。最终模型在 DS 和 RF 接口下分别评估，记为
% RePaIR\(_{\mathrm{DS}}\) 和 RePaIR\(_{\mathrm{RF}}\)。

% ReaSon 与 RS 的比较用于分析 rationale 来源的作用，PaIR 与
% Mix 的比较用于分析分支级归一化的作用，RePaIR 与 PaIR
% 的比较则用于分析替换 rationale 来源带来的增量影响。

\subsection{Regenerated-Rationale Paired-Interface Reweighting (RePaIR)}

RePaIR combines the rationale replacement of ReaSon with the paired
supervision of PaIR. Starting from the base model \(\theta_0\), the
first stage trains a PaIR model \(\theta_A\) on the teacher-rationale
dataset \(\mathcal{D}_T\). The resulting model then regenerates
rationales \(r_i^A\) under RF. These rationales are paired with the
gold scores \(s_i\) to form
\(\mathcal{D}_A=\{(x_i,r_i^A,s_i)\}\).

The second stage reinitializes the model from \(\theta_0\) and applies
PaIR to \(\mathcal{D}_A\), rather than further updating
\(\theta_A\). The resulting model is evaluated under DS and RF,
denoted by RePaIR\(_{\mathrm{DS}}\) and
RePaIR\(_{\mathrm{RF}}\), respectively.

Comparing ReaSon with RS assesses the effect of rationale source.
Comparing PaIR with Mix assesses branch-wise normalization.
Comparing RePaIR with PaIR assesses the additional effect of
replacing teacher rationales with student-generated rationales.
